\section{A neutral-current obstruction to absolute extensivity}
\label{sec:counterexamples}

\begin{theorem}[Superextensive vacuum trial energy]
\label{b-vacuum-L5}
Fix $m,\Lambda_f,\Lambda_\gamma,\alpha>0$ in the bare model with
massless photons. There are constants $c_\Lambda,C_\Lambda>0$ and $L_0<\infty$
such that for all $L\ge L_0$ its exact neutral interacting vacuum obeys
\begin{equation}
 E^{\vac}_{\Lambda,L}\le-c_\Lambda L^5+C_\Lambda L^3.
 \label{eq:vacuum-instability}
\end{equation}
The same conclusion holds for density-product ordering. In either
case no absolute bound of the form \eqref{eq:W} holds in this fixed
normalization, even with no nuclei.
\end{theorem}
\begin{proof}
Set $q=2\pi/L$ and let $C_L=P_f\cos(qx_1)P_f$ on the scalar band.
Its sign $S_L=\sign C_L$, with $\sign(0)=1$, is a self-adjoint
unitary. The one-particle matrix
\begin{equation}
 \gamma_L=\tfrac12(I+\mathsf a_2S_L)
 \label{b-L5-projector}
\end{equation}
is a projection of rank $2M_L$: the spin trace of $\mathsf a_2$
is zero. Its Slater determinant therefore has exactly zero charge.
Moreover $\tr\gamma_L(x,x)=2d_L=\tr P_-(x,x)$, so its charge
density vanishes pointwise. Its transverse current is
\[
 J_L=\int\cos(qx_1)\langle j_2(x)\rangle\dd x
 =\Tr[(\gamma_L-P_-)\mathsf a_2C_L]=2\Tr|C_L|.
\]
The sea contribution is zero because $C_L$ has zero momentum diagonal.
As $\|C_L\|\le1$, counting the two momentum transfers gives
\[
 J_L\ge2\Tr C_L^2
 =\#\{p\in\Gamma_{f,L}:p+qe_1\in\Gamma_{f,L}\}
 \ge c_{\Lambda_f}V
\]
for all sufficiently large $L$. For example, count an interior ball
of fixed positive radius inside the cutoff ball.

The kinetic energy is at most $4M_LE_{\Lambda_f}=O(V)$.
For clarity, the full Wick Coulomb expectation of a quasi-free
state with $Q_L=\gamma_L-P_-$ is exactly
\begin{equation}
 \langle W_P(T_{-k},T_k)\rangle
 =\Tr(Q_LT_{-k})\Tr(Q_LT_k)
       -\Tr(Q_LT_{-k}Q_LT_k).
 \label{eq:trial-wick}
\end{equation}
This follows by the CAR Wick contractions in its density matrix.
The direct term is zero by the pointwise charge identity, whereas
$|\Tr(Q_LT_{-k}Q_LT_k)|\le\|Q_L\|_{\HS}^2\le4M_L$.
Equation~\eqref{eq:coulomb-mode-sum} therefore bounds the complete
Coulomb expectation in absolute value by $C_\Lambda V$.
This includes the charge fluctuations; they have not been set to zero.

Choose a photon coherent state with mean field
$A(x)=a e_2\cos(qx_1)$. It is transverse and its nonzero modes
are retained for all large $L$. The field energy and coupling are
\[
 \frac{Vq^2a^2}{4}+eaJ_L.
\]
Minimizing in the real amplitude gives
$-e^2J_L^2/(Vq^2)\le-c_\Lambda L^5$.
Adding the preceding $O(V)$ costs proves
\eqref{eq:vacuum-instability}. Every trial state belongs to the
exact neutral sector and has finite photon occupation at its fixed
volume. The ordering difference is bounded by $B_\Lambda V$ from
\eqref{eq:full-coulomb-bound}, so the same proof applies to
density-product ordering.
\end{proof}

The optimized mean vector potential is $O(L^2)$ and its magnetic
field is $O(L)$. A momentum cutoff alone imposes no bound on these
amplitudes. A declaration that the model describes low energies
cannot exclude these allowed trial states without changing the
stated space or Hamiltonian. Fixed mass counterterms and a fixed
vacuum-energy density contribute at most $O(V)$ here. A fixed
nonzero charge renormalization changes the coefficient, not the
power of $L$.

There is no contradiction with \eqref{b-headline}: the latter
subtracts the actual interacting vacuum and compares the change
caused by nuclei. A large negative common vacuum contribution does
not disprove that relative inequality. Conversely, separate bounds
$E(R,Z)\ge-AK-BV$ and $-BV\le E^{\vac}\le0$ do not imply a
relative bound $-CK$; they leave a volume error. The canonical
Gaussian-domination argument supplies the comparison that these
separate estimates lack.

The ordinary absolute vacuum-energy density for this bare regulator
with massless photons cannot have a finite thermodynamic limit: the trial bound
already gives $E^{\vac}/V\to-\infty$ in upper bound. Subtracting
the exact finite-volume vacuum is a legitimate relative comparison,
but does not by itself construct a thermodynamic state or establish
physical renormalization.
